\section{问题三的模型建立与求解}

\subsection{任务条件与时间账本}

问题三的全部干扰源均为全向源，位置静止，频道唯一，接收半径未知但取值于 $[1000,1500]$ 米，总数 $N$ 未知且 $10\le N\le16$。机器狗自原点出发，初始检测频道为 1，移动速度 5 米/秒，清除操作不改变检测频道，也不要求返回原点。

在确保全部干扰源被清除的前提下，本文以虚拟总时间 $T$ 为优化目标。由式 \eqref{eq:time-ledger}，整场耗时可分解为移动、检测、换频与清除四项之和。该分解的意义在于把"少走路、少检测、少切频道、少失败尝试"放在同一时间尺度上，使局部决策不会因只优化其中一项而损害整体。需要说明的是，机器狗自身的计算时间不计入虚拟时间，但受题目给出的程序运行时间上限约束。

由于不同场景的目标数 $N$ 不同，跨场景对照一律采用\textbf{同一场景配对}；主指标为场景等权的均值 $K^{-1}\sum_j T_j/N_j$。本文另存合并比值 $\sum_jT_j/\sum_jN_j$ 用于核对，但不与场景等权均值混用。评测时不删除最后一次成功清除之后的必要搜索时间——在程序尚不知道 $N$ 的情况下，这段收尾是完成"确认无遗漏"所必需的。

\subsection{单源定位区域与保证清除}

对每个频道维护其全部正向测量历史，并求取凸外包多边形 $P$。先验取半径 1800 米圆域的外接正 180 边形；每次正向测量与半宽 $\varepsilon$ 的正向角域、以及以测站为中心半径 1500 米圆盘的外接多边形相交；近距（过强信号）记录与半径 5 米圆盘的外接多边形相交。角域半平面与最小包围圆的代码逐字复用问题一的实现并记录哈希。

角域半宽取 $\varepsilon=1.0051^\circ$，包含 $\pm1^\circ$ 的物理误差、读数保留两位小数带来的 0.005$^\circ$ 舍入以及 0.0001$^\circ$ 的数值余量。由假设五，同一位置的误差固定，不能通过在同一点重复测量来缩小误差。本文忽略近距记录对应的半径 5 米内部排除区，也忽略无信号记录产生的非凸排除区，这属于保守放宽，而不是把"无信号"当作"无目标"。

设 $P$ 的顶点为 $v_j$，最小包围圆为 $(c,\varrho)$。仅当 $\varrho\le19.5$ 米时才尝试有证书的清除；具体清除点 $q$ 满足
\begin{equation}
  \max_j\lVert q-v_j\rVert\le19.5 .
  \label{eq:q3-clear}
\end{equation}
对凸多边形，范数的最大值可在顶点取得，因此真实源到 $q$ 的距离不超过该上界，即不大于 19.5 米，留出 0.5 米的数值与工程余量。与问题一一致，$D(P)\le40$ 不能替代这一条件。

为节省路程，本文在"圆心 $c$ 指向机器狗当前位置 $p$"的线段上，用 45 次二分求最靠近 $p$ 的安全点；若 $p$ 已满足约束则原地清除。这是该线段上的最优点，不声称是全部安全清除点集到 $p$ 的全局投影。每次清除前保存顶点、点位与最坏距离证书；若接口返回失败，则保留冲突并退出，不得悄悄忽略。

\subsection{负观测的信息与接收半径约束}

无信号记录在第三问中携带两类信息。

\textbf{第一类：圆盘排除。} 若在位置 $n$ 未接收到某频道信号，则该源到 $n$ 的距离超过其接收半径，而接收半径不小于 1000 米，故可排除以 $n$ 为中心、半径 1000 米的圆盘。该操作通常产生非凸集；实现中使用
\begin{equation}
  \operatorname{conv}\Bigl(P\setminus\bigcup_{n\in S^-}\overline{B}(n,999.8)\Bigr)
  \label{eq:disk-exclusion}
\end{equation}
作为安全的凸外包，不把"圆外"误写成某个方向上的半平面。候选顶点由保留的多边形顶点、边与圆的交点、圆与圆的交点构成，过滤掉落在其他排除圆内或 $P$ 外的点后再求凸包。使用 999.8 米而非 1000 米提供数值余量。由于圆盘外侧边界对可行集而言是向内凹的，线性函数的最大值无需取其弧内部点。

\textbf{第二类：正负观测对生成的半平面。} 设 $s$ 为任一次有信号观测位置。由于同一个源的接收半径 $R$ 固定，有
\begin{equation}
  \lVert g-s\rVert\le R<\lVert g-n\rVert
  \;\Longrightarrow\;
  2(n-s)^{T}g<\lVert n\rVert^2-\lVert s\rVert^2 .
  \label{eq:pos-neg-halfplane}
\end{equation}
将每一对正、负观测生成的半平面加入 $P$，并使用向外 $10^{-6}$ 米的线性边界余量。该结论不要求接收半径恰为 1000 米，也不适用于接收半径随时间变化的模型。

\subsection{连续场地的停止证据}
\label{sec:coverage}

这是问题三中最关键、也最容易出错的部分。题目说明源总数在 10 至 16 之间但具体个数未知，因此\textbf{清除了 10 个或 14 个都不能结束}。本文只允许两种终止依据。

\textbf{依据一：}已清除 16 个源，达到题目给出的总数上界。

\textbf{依据二：}全部已发现源均已清除，并且每一个未发现频道都有全场覆盖证书。

覆盖证书的构造如下。将 $[-1800,1800]^2$ 等分为 $128\times128$ 个边长为 28.125 米的方格，保留所有与目标圆域相交的方格（含部分落在圆外的边界格）。对某个频道的一次无信号记录，只有当该方格\textbf{四个角点}到该无信号位置的距离都不超过 999.99 米时，才把整个方格标记为排除；由圆盘的凸性，这保证整格中的每一点都确实被排除。若某频道的全部保留方格均被其无信号记录覆盖，则该频道不可能存在源。

这一判据是\textbf{整格距离证书}，而不是仅检查网格中心的采样，两者有本质区别：后者在方格尺度较大时会漏掉真实的可能位置。

作为冗余核查，7 个外围测站的覆盖还可以直接验证：任一点与某个外围站的极角差不超过 $30^\circ$，比较其到原点及该外围站的距离，沿半径分段后最大最近站距离出现在半径 1800 米、角差 $30^\circ$ 的边界点处，约为 924.9 米，小于 1000 米。代码进一步对所有保留方格的完整覆盖进行独立检查。

此外还需注意两点。其一，发现 16 个不同频道后可以停止继续搜索未知频道，但仍必须清除这 16 个已发现源。其二，算法中"计划在某点扫描"只影响下一步的路线评分，\textbf{不写入实际排除记录}——未来覆盖不能当作已完成覆盖。

\subsection{发现—定位—清除闭环与联合路线}

\textbf{初始扫描。} 先在原点扫描所有尚未被排除的频道并缓存全部测量。同一位置的 20 个频道分别计时。

\textbf{任务点集合。} 把已知目标的安全清除点（或区域中心）作为必访点，把足以覆盖剩余场地的扫描点作为附加点。用固定起点、自由终点的最近邻路径初始化，再做 2-opt 改善顺序；每次实际观测后重新规划，只执行当前路线的第一项。因此这是一个滚动时域启发式，而不是全局最优动态规划。

\textbf{定位补测。} 初始设置一段 250 米的短基线用于同时改善多个已知目标的交会角。对已知频道的补测要求：距旧测点至少 150 米，区域尚不满足清除条件，且预计距离不超过 1200 米。对已有两次以上测量的目标，用中心方位及 $\pm1^\circ$ 三种假想读数预览区域收缩程度。需要强调，这个预览是补测价值的启发式估计，三种读数不是连续误差范围的最坏情况证书；真实清除仍使用完整的有界区域检查。

\textbf{扫描点的连续调整。} 从"目标中心加六个外圈站"构成的可行覆盖出发，逐一删除冗余扫描点。删除评分折算为移动距离：约 $30U$ 米（$U$ 为未发现频道数，每次检测近似 6 秒，速度 5 米/秒）的未知频道检测成本，加上可能缩短的局部路线。固定其余扫描站，对第 $i$ 站独占负责的剩余方格，取其全部角点的凸包顶点集 $V_i$，令前后路线点为 $a$、$b$，求解二维凸子问题
\begin{equation}
  \min_q\ \lVert q-a\rVert+\lVert q-b\rVert
  \qquad\text{s.t.}\quad \lVert q-v\rVert\le999.99\quad(\forall v\in V_i).
  \label{eq:scan-relocate}
\end{equation}
末站省略到 $b$ 的距离。将坐标缩放至千米后用 SLSQP 数值求解；只接受所有原始距离约束通过且局部目标不增加的结果，随后重新检查整体网格覆盖。不以求解器的成功标记代替可行性检查。删除、移动与路径重排组成交替改进；整体组合问题不是凸问题，本文不宣称全局最优。

\textbf{递推式的覆盖思路。} 上述覆盖构造借鉴了"首步可行 $\to$ 相邻步保持覆盖 $\to$ 末步覆盖边界 $\to$ 停止"的递推论证顺序，使任务完成有明确的判据，而不是"运行到次数上限"。但必须强调：连续测线在旧题中天然覆盖带内全部点，而本题中\textbf{走过一段路不等于在这段路上测过全部频道}；接收半径未知时，保证性论证只能使用已知下界 1000 米，不能使用 1500 米或样本平均半径。

\subsection{光学动作的时间价值}

在定位区域已经足够小时，直接进行一次光学试探可能比继续测向更省时间。本文仅在 $19.5<\varrho\le100$ 米时比较试清动作。候选位置包括当前位置、包围圆中心、朝当前机器狗偏移的点以及长轴两侧的点，每源最多试探两次；位置信息表示点来自当前的非凸可行集 $F$ 而非模拟器真值。

\textbf{非凸失败信息的保留。} 已知源在 $q$ 点清除失败意味着 $\lVert g-q\rVert>20$，这是一条有效信息。本文在数值实现中移除"半径 19.999 米圆盘的内接 64 边形"，所得删除区域严格包含在真实可排除范围内，因而留下的仍是外包集合。实现保持两种表示：非凸集合 $F$ 保存全部光学失败排除，$\operatorname{conv}(F)$ 的顶点用于原有测向、接收与最小包围圆计算；后续正向测向区域与 $F$ 相交，不以凸包覆盖 $F$ 的存储。保留所有连通分量与孔洞。

\textbf{动作评价。} 对动作 $a$ 与代表情景 $h$，评价量取
\begin{equation}
  C(a,h)=\frac{\text{移动距离}}{5}+\text{本动作检测或清除费用}+\text{后续常规定位完成费用},
  \label{eq:action-value}
\end{equation}
其中试清成功计 5 秒、失败计 3 秒并使用更新后的区域预测后续费用。无线电预测每次按 6 秒保守计入切换，而实际模拟按真实频道计费。预测使用至多 5 步原几何定位；评价取 $0.8$ 倍情景均值加 $0.2$ 倍情景最大值，且至少比无线电方案少 1 秒才尝试试探。需要说明：均值权重与代表情景都是规划约定，不是概率模型，也不是连续最坏情形的证明；预测并未完整建模试清后对整场路径的影响，因此最终必须经实际全场验证。这一"先看本轮动作如何影响未来，再向后试算若干步"的组织方式，参考了往年优秀论文中有限步状态更新与"当前成本加后续成本"的论证结构，但本题的物理条件与计费规则均按今年附件重新建立，原论文中的免费多目标观测、概率噪声与全局真值假设均未沿用。

\subsection{求解结果}

\subsubsection{同场配对审计}

本文使用固定随机种子的合成场景集进行同场配对比较，所有结果均计入整场末尾排除与失败清除费用。

\begin{table}[H]
  \centering
  \caption{问题三的同场配对审计结果（场景等权 $T/N$，单位：秒/点）}
  \label{tab:q3-results}
  \begingroup
  \renewcommand{\arraystretch}{1.4}
  \begin{tabularx}{\textwidth}{L{0.34\textwidth}C{0.11\textwidth}C{0.14\textwidth}C{0.14\textwidth}C{0.13\textwidth}}
    \toprule[1.2pt]
    {\heifont\bfseries 测试组} & {\heifont\bfseries 场数} & {\heifont\bfseries 对照版} & {\heifont\bfseries 本模型} & {\heifont\bfseries 平均改善} \\
    \midrule[0.5pt]
    200 场同场审计（联合规划 vs 基线） & 200 & 317.38 & 237.19 & 25.27\% \\
    100 场三版本对照（v3 / v4 / v5） & 100 & 238.11 & 234.23 & 1.63\% \\
    60 场冻结配对（v5 $\to$ v6 光学评价） & 60 & 234.03 & 232.60 & 0.61\% \\
    \bottomrule[1.2pt]
  \end{tabularx}
  \endgroup
\end{table}

在 200 场同场配对审计中（种子 5100--5299），基线、区域耦合、连续覆盖与组合方案的场景等权均值依次为 317.38、238.71、238.40、237.19 秒/点；四套方案全部实现 200/200 完整清除、零失败清除，配对均值的近似 95\% 区间为 $[232.52,241.85]$ 秒/点。需要说明：该区间只表示给定本地生成规则下的场景抽样不确定性，\textbf{不是}官方泛化区间。

在 100 场三版本同场对照中，v3、v4、v5 分别为 238.1110、234.8970、234.2317 秒/点，全部完整清除、零失败清除。v4 与 v5 的配对均值差为 0.6652 秒/点，近似 95\% 区间为 $[-1.8046,3.1351]$，包含 0；因此\textbf{不能认定 v5 相对 v4 具有稳定优势}。这一轮实验中平均每场少走 322.15 米但多检测 8.93 次，测量费用抵消了大部分移动收益，这为"补测价值估计"这一后续方向提供了证据。

在 60 场冻结留出配对中（种子 17000--17059），加入光学动作评价的 v6 相对 v5 由 234.03 降至 232.60 秒/点，平均节省 1.4283 秒/点（0.61\%），配对近似 95\% 区间为 $[1.0976,1.7591]$，54 场改善、6 场退化，最差场景（种子 17020）增加 1.7588 秒/点。本轮平均路程由 10833.10 米降至 10816.56 米，检测次数由 123.28 降至 118.73 次，新版平均发生 3.33 次失败光学尝试，其时间代价均已计入。

\begin{figure}[H]
  \centering
  \includegraphics[width=0.94\textwidth]{q3-joint-paired.png}
  \caption{问题三的逐场配对比较。200 场同场审计中各方案相对基线的每场时间变化与累计分布，含末尾排除与失败清除费用}
  \label{fig:q3-joint}
\end{figure}

\subsubsection{真实官方演练}

在本地合成实验之外，本策略完成了一次真实的官方演练，案例编码 JCEK-BMVS-WD8Z-WF6J。官方结果文件确认场内共有 14 个全向干扰源；动作响应确认成功清除 14 个、零失败，正常主动退出。总虚拟时间 3854.083567 秒，平均每源 275.29 秒，累计移动 16020.418 米。全部数据来自保存的官方接口响应，并已按原始动作独立重算；未解密官方日志，也未连接或重跑模拟器。

\begin{table}[H]
  \centering
  \caption{问题三官方演练的时间分解（案例 JCEK-BMVS-WD8Z-WF6J，14 个全向源）}
  \label{tab:q3-official}
  \begingroup
  \renewcommand{\arraystretch}{1.4}
  \begin{tabularx}{\textwidth}{L{0.34\textwidth}C{0.18\textwidth}C{0.16\textwidth}Y}
    \toprule[1.2pt]
    {\heifont\bfseries 计费项} & {\heifont\bfseries 虚拟时间/秒} & {\heifont\bfseries 占比} & {\heifont\bfseries 次数} \\
    \midrule[0.5pt]
    移动            & 3204.084 & 83.13\% & 累计 16020.42 米 \\
    检测            & 495.000  & 12.84\% & 99 次 \\
    检测频道切换    & 85.000   & 2.21\%  & 85 次 \\
    成功清除        & 70.000   & 1.82\%  & 14 次 \\
    清除失败        & 0.000    & 0.00\%  & 0 次 \\
    \midrule[0.5pt]
    合计            & 3854.084 & 100.00\% & 115 个唯一动作 \\
    \bottomrule[1.2pt]
  \end{tabularx}
  \endgroup
\end{table}

本局共发出 115 个唯一动作（进入 1、检测 99、清除 14、退出 1），115 个请求与 115 个响应全部一次接受，无重试、拒绝或未确认动作；逐动作时间守恒的最大偏差为 $4.87\times10^{-7}$ 秒，符合响应的小数舍入。39 次正向观测中，14 次为首次发现，其余 25 次为定位补测，即每个源平均新增 1.79 次定位检测。按动作用途划分，扫描点检测 74 次、专门定位检测 25 次。最后一次清除发生在 3292.188 秒，之后用去 561.896 秒（占全程 14.58\%）完成余下两个测站、12 次无信号检测。

这一收尾段不是可以直接删除的浪费：程序当时并不知道总数为 14，仍需排除其他频道；第 4 个测站虽然没有发现新源，但其无信号记录参与了完整覆盖证书。这正说明了第 \ref{sec:coverage} 节终止依据的必要性。

\begin{figure}[H]
  \centering
  \includegraphics[width=0.88\textwidth]{q3-official-route.png}
  \caption{问题三官方演练的实际轨迹与清除进度。绿色标记为清除动作发生位置，其中部分为原地近距清除；图中标记是清除点而非干扰源真实坐标}
  \label{fig:q3-official-route}
\end{figure}

图 \ref{fig:q3-official-route} 给出该局的轨迹与逐源清除进度。需要特别指出，图中的清除标记是机器狗执行清除时的位置（距真实源不超过 20 米），\textbf{不能当作干扰源的真实坐标}。

\subsubsection{与前次演练的比较及可比性说明}

本策略另与两次更早的演练记录作了对比，结果见表 \ref{tab:q3-history}。

\begin{table}[H]
  \centering
  \caption{问题三三次官方演练记录对比}
  \label{tab:q3-history}
  \begingroup
  \renewcommand{\arraystretch}{1.4}
  \begin{tabularx}{\textwidth}{L{0.26\textwidth}C{0.10\textwidth}C{0.13\textwidth}C{0.12\textwidth}C{0.10\textwidth}C{0.10\textwidth}}
    \toprule[1.2pt]
    {\heifont\bfseries 演练} & {\heifont\bfseries 已清除} & {\heifont\bfseries 总虚拟时间/秒} & {\heifont\bfseries 每源平均/秒} & {\heifont\bfseries 路程/km} & {\heifont\bfseries 清除失败} \\
    \midrule[0.5pt]
    20260910-200250 & 14 & 6446.68 & 460.48 & 27.248 & 11 \\
    20260910-222119 & 16 & 6684.96 & 417.81 & 27.595 & 15 \\
    20260911-002124（本策略） & 14 & 3854.08 & 275.29 & 16.020 & 0 \\
    \bottomrule[1.2pt]
  \end{tabularx}
  \endgroup
\end{table}

相对第一局 14 源旧演练，本局总时间低 40.22\%、路程少 41.21\%；相对第二局 16 源旧演练，每源平均时间低 34.11\%。\textbf{但这三局的案例并不相同}，目标数相同也不代表场景相同，因此这些差值不能作为策略因果提速的证据，也不能替代同案例配对比较。本文将其列为实际运行记录，仅用于说明策略在真实接口上可完整运行。

\subsection{有效性检验与未达标说明}

\textbf{几何与动作的独立核验。} 对演练日志，本文重建了 39 次正向观测的完整约束，对每个区域做 16 方向线性规划支持函数检查共 624 项，最大差异 $1.23\times10^{-11}$ 米；14 次清除点全部通过顶点最坏距离复核。对未发现的频道 2、5、6、7、19、20，均有 7 个测站的真实无信号记录，从动作记录重新构造的整格覆盖最小半径余量为 60.621 米——这一结论不是仅相信摘要中的完成标志。

\textbf{压力与消融测试。} 第三问的边界布局、固定极端测角误差（恒定 $+1^\circ$ 与 $-1^\circ$）以及棋盘式误差场测试均实现完整清除，但极端误差组略慢、棋盘误差组几乎持平（-0.00\%）。20 场消融实验中，"不更新失败区域"为 238.0196 秒/点，"更新失败区域"为 237.9988 秒/点，差值仅 0.0208，区间 $[-0.3224,0.3640]$；因此\textbf{不能把整体提速归因于失败信息更新这一单项}——它在几何上是有效的，但效率贡献尚未被证实。

\textbf{200 秒/点目标未达成。} 需要如实说明：本轮全部讨论中出现的"200 秒/点"是团队设定的内部研究目标，并非题目要求。当前方案未达到该目标。进一步地，本文构造的圆周 10 源场景给出移动下界 233.64 秒/点，这否定了"在全部合法场景上获得 200 秒硬保证"的可能性；但这\textbf{不}否定特定分布下均值可能降至 200 秒，本文也没有证明该分布均值不可达。

此外，16 个源子组的平均 193.09 秒/点不能替代全体场景的 237.19 秒/点。不能选择容易场景、忽略失败或删去收尾段后宣称达成目标。

\textbf{剩余的主要研究方向}是跨源定位任务、停靠位置与未来路径的联合选择。本轮的局部完成费用预测不能称为完整的多步全局规划，独立审计也显示当前局部成本的改善并没有稳定减少整场移动距离。

\subsection{正式测试结果}

按题目要求，需在充分演练后进行三次正式测试，并按表 1 的格式给出结果。

\begin{table}[H]
  \centering
  \caption{问题三正式测试结果}
  \label{tab:q3-formal}
  \begingroup
  \renewcommand{\arraystretch}{1.45}
  \begin{tabularx}{\textwidth}{L{0.22\textwidth}C{0.18\textwidth}C{0.20\textwidth}C{0.20\textwidth}}
    \toprule[1.2pt]
    {\heifont\bfseries 测试案例编码} & {\heifont\bfseries 清除干扰源个数} & {\heifont\bfseries 平均定位清除时间} & {\heifont\bfseries 程序运行时间} \\
    \midrule[0.5pt]
    测试1 \hspace{2em} & \hspace{2em} & \hspace{2em} & \hspace{2em} \\
    测试2 \hspace{2em} & \hspace{2em} & \hspace{2em} & \hspace{2em} \\
    测试3 \hspace{2em} & \hspace{2em} & \hspace{2em} & \hspace{2em} \\
    \bottomrule[1.2pt]
  \end{tabularx}
  \endgroup
\end{table}

\textbf{截至本文写作时，三次正式测试尚未进行，表 \ref{tab:q3-formal} 的内容留空。} 本文已有的全部第三问数值均来自本地合成场景与一次官方演练，不构成正式测试成绩。相应地，正式测试对应的三份模拟器日志文件也将在测试完成后导出并放入支撑材料。
